\section{引言：从医学誓言到 AI 伦理}

本讲是 MIT 6.S191（Introduction to Deep Learning）2025 年春季课程的第 9 讲，由 Comet ML 研究负责人 Doug Blank 作为嘉宾主讲。Doug Blank 自 1990 年代起从事神经网络研究，1997 年完成认知科学与计算机科学博士学位，曾在 Bryn Mawr College 担任教授 20 年，后加入 AI 实验管理初创公司 Comet ML。

本讲的核心问题极具思辨性：\textbf{我们能否让 AI 系统做出类似医生的``希波克拉底誓言''（Hippocratic Oath）——承诺``不造成伤害''（Do No Harm）？}

\begin{importantbox}{希波克拉底誓言（Hippocratic Oath）}
希波克拉底誓言是约 2500 年前由古希腊医学之父 Hippocrates 提出的医学伦理准则。医学生在正式执业前宣誓遵守专业伦理标准，核心原则可概括为一句话：\textbf{Do No Harm}（不造成伤害）。它体现了两个关键原则：
\begin{itemize}
    \item \textbf{Beneficence}（行善原则）：积极为患者谋求利益
    \item \textbf{Non-maleficence}（不伤害原则）：避免对患者造成损害
\end{itemize}
\end{importantbox}

讲者指出，深度学习系统面临的问题远不止技术层面——训练数据质量、过拟合、泛化能力等——还包括数据偏见、能源消耗、输出准确性、安全性等更广泛的社会问题。随着系统规模和复杂度的增长，这些问题也在不断加剧。

\begin{knowledgebox}{从 1997 到 2025：深度学习的规模跃迁}
讲者以自身经历说明了深度学习的规模变化：1997 年的博士论文使用的神经网络仅有约 30 个神经元和 3000 个参数，而如今的大语言模型拥有数千亿参数。数学原理基本相同——矩阵乘法、激活函数、反向传播——但硬件效率的提升使得系统规模呈指数级增长，随之而来的是全新的工程和伦理挑战。
\end{knowledgebox}

\subsection{本章小结}

本讲围绕一个核心类比展开：如果医学从业者需要宣誓``不造成伤害''，那么 AI 从业者是否也应该做出类似的承诺？讲者从投资学中的风险-收益框架切入，逐步探讨了 AI 项目的风险评估、法规监管、技术局限和开发者责任。

